\documentclass[10pt,a4paper]{article}

% Template de relatório — Trabalho de Visão Computacional
% CIFAR-10: CNN própria e Transfer Learning
% INSTRUÇÕES: consulte os comentários ao longo deste arquivo para saber o que
% deve ser incluído em cada seção. Os comentários não aparecem no PDF compilado.
% Remova os textos entre colchetes antes da entrega final.
% As tabelas são exemplos e devem ser adaptadas aos experimentos realizados.
% Limite recomendado: 10 páginas, incluindo referências.

\usepackage[utf8]{inputenc}
\usepackage[T1]{fontenc}
\usepackage[brazil]{babel}
\usepackage{lmodern}
\usepackage[a4paper,margin=2.2cm]{geometry}
\usepackage{microtype}
\usepackage{graphicx}
\usepackage{booktabs}
\usepackage{multirow}
\usepackage{amsmath,amssymb}
\usepackage{url}
\usepackage{xcolor}
\usepackage{caption}
\usepackage{subcaption}
\usepackage{float}
\usepackage{enumitem}
\usepackage{hyperref}

\hypersetup{
    colorlinks=true,
    linkcolor=blue,
    citecolor=blue,
    urlcolor=blue
}

\setlength{\parindent}{1.25cm}
\setlength{\parskip}{0pt}
\emergencystretch=2em

% Substitua os campos de identificação.
\title{\vspace{-1.2cm}\Large\bfseries [Título do artigo]}
\author{
    [Nome do integrante 1]\textsuperscript{1},
    [Nome do integrante 2]\textsuperscript{1},
    [Nome do integrante 3]\textsuperscript{1}\\[3pt]
    \small \textsuperscript{1}[Departamento, instituição]\\
    \small \texttt{[e-mails dos integrantes]}
}
\date{}

\begin{document}
\maketitle
\vspace{-0.8cm}

% Escreva um resumo conciso do problema, objetivo, métodos,
% protocolo experimental, principais resultados e conclusão.
\begin{abstract}
[Resumo: problema e objetivo; abordagens avaliadas; principais resultados quantitativos; conclusão principal.]
\end{abstract}

\noindent\textbf{Palavras-chave:} [3 a 5 palavras-chave.]

% ============================================================
\section{Introdução}
% Instruções da especificação:
% Apresente o problema de classificação de imagens, o dataset CIFAR-10 e a
% motivação para utilização de redes neurais convolucionais e transfer learning.
% Ao final da introdução, apresente claramente os objetivos do trabalho e,
% preferencialmente, as principais questões que serão investigadas experimentalmente.

Classificar imagens é uma tarefa fundamental em visão computacional, mas a
variação de pose, iluminação, fundo e aparência entre objetos torna difícil
definir regras manuais que funcionem para todas as imagens. Redes neurais
convolucionais aprendem representações diretamente dos pixels e exploram a
estrutura espacial por meio de filtros aplicados localmente e compartilhados
ao longo da imagem \cite{lecun1998}. O CIFAR-10 é um conjunto de referência
para esse problema, composto por imagens coloridas de baixa resolução
distribuídas em dez classes \cite{krizhevsky2009}.

Além de desenvolver uma CNN própria, o trabalho investiga modificações
arquiteturais e o uso de modelos pré-treinados. No transfer learning, parte das
representações aprendidas em uma tarefa de origem pode ser adaptada à tarefa
alvo; a utilidade dessas representações depende, entre outros fatores, da
camada transferida e da relação entre as tarefas \cite{yosinski2014}. Essa
abordagem será comparada com o treinamento da CNN própria sob um protocolo
comum.

O objetivo é comparar o desempenho das arquiteturas e das estratégias de
treinamento no CIFAR-10, considerando também o custo em parâmetros e tempo e
as diferenças de desempenho entre classes. As configurações serão selecionadas
com os dados de validação. O conjunto de teste permanecerá separado durante o
desenvolvimento e será utilizado somente na avaliação final das configurações
selecionadas.

% ============================================================
\section{Materiais e métodos}

\subsection{Conjunto de dados e protocolo experimental}
% Instruções da especificação:
% Descreva o dataset CIFAR-10, a divisão dos dados utilizada, o pré-processamento,
% as técnicas de data augmentation (quando utilizadas) e o protocolo experimental.
% Informe a divisão de 45.000 imagens para treinamento, 5.000 para validação e
% 10.000 para teste, explicando que a divisão treino/validação é estratificada,
% fixa para o grupo e mantida em todos os experimentos.

O CIFAR-10 contém 60.000 imagens coloridas de dimensão $32\times32$, divididas
em dez classes. O conjunto oficial de treinamento contém 50.000 imagens e o
conjunto de teste contém 10.000 imagens \cite{krizhevsky2009}. Para o
desenvolvimento, os 50.000 exemplos de treinamento foram divididos uma única
vez, de forma estratificada, em 45.000 exemplos para treinamento e 5.000 para
validação, com semente fixa igual a 42. A validação contém 500 exemplos de cada
classe. Os mesmos índices são reutilizados em todos os experimentos do grupo;
os exemplos oficiais de teste não são usados na seleção de modelos ou
hiperparâmetros.

No pré-processamento padrão, as imagens são convertidas para tensores e
normalizadas por canal usando média $(0.4914, 0.4822, 0.4465)$ e desvio-padrão
$(0.2470, 0.2435, 0.2616)$. Quando habilitada, a ampliação de dados é aplicada
somente ao treinamento e consiste em espelhamento horizontal aleatório e
recorte aleatório de $32\times32$ após preenchimento de quatro pixels. A
validação não recebe transformações aleatórias. A infraestrutura permite
redimensionamento e normalização ImageNet para os experimentos que os
requeiram; os valores efetivamente usados por cada modelo devem ser registrados
junto aos respectivos hiperparâmetros.

O pipeline fixa as sementes do experimento antes de carregar os dados e
embaralha apenas os lotes de treinamento. A perda, o otimizador, a taxa de
aprendizado, o tamanho do lote, o número de épocas e o scheduler são definidos
na configuração de cada execução. A seleção de checkpoints e ajustes é feita
pela validação. Os hiperparâmetros finais da CNN própria e de cada modelo
pré-treinado são apresentados nas subseções correspondentes.

\subsection{CNN própria e modelo baseline}
% Instruções da especificação:
% Descreva a arquitetura baseline desenvolvida, incluindo número de camadas,
% filtros, kernels, funções de ativação, pooling, camadas totalmente conectadas,
% Batch Normalization e Dropout (quando aplicáveis) e número total de parâmetros.
% Descreva também a função de perda, o otimizador, os hiperparâmetros e as
% decisões tomadas para o treinamento. Justifique as escolhas sempre que possível.

[Descrição da arquitetura baseline e dos hiperparâmetros.]

% TABELA DE EXEMPLO: adapte a estrutura à arquitetura efetivamente implementada;
% adicione, remova ou altere linhas e colunas conforme necessário.
\begin{table}[H]
\centering
\caption{Exemplo de organização da arquitetura baseline da CNN própria.}
\label{tab:baseline}
\small
\begin{tabular}{llll}
\toprule
Camada & Configuração & Dimensão de saída & Parâmetros\\
\midrule
Camada 1 & [Configuração] & [Dimensão] & [Número]\\
Camada 2 & [Configuração] & [Dimensão] & [Número]\\
Adicionar linhas conforme necessário & & &\\
\bottomrule
\end{tabular}
\end{table}

\subsection{Experimentos de modificação arquitetural}
% Instruções da especificação:
% Descreva as modificações realizadas na arquitetura e justifique as escolhas.
% Informe como os experimentos foram organizados e quais fatores foram mantidos
% constantes. Se algum hiperparâmetro foi ajustado em função de uma alteração
% arquitetural, justifique e documente o ajuste. Discuta como cada alteração
% permite investigar seu impacto no desempenho.

[Descrição e justificativa dos experimentos arquiteturais.]

% TABELA DE EXEMPLO: os experimentos listados são ilustrativos; substitua-os pelas
% modificações realmente investigadas e ajuste a tabela ao seu protocolo.
\begin{table}[H]
\centering
\caption{Exemplo de organização dos experimentos de modificação arquitetural.}
\label{tab:experimentos}
\small
\begin{tabular}{clp{7cm}}
\toprule
ID & Configuração & Modificação em relação à baseline\\
\midrule
B0 & Baseline & Arquitetura de referência\\
E1 & [Nome] & [Descrição da alteração]\\
E2 & [Nome] & [Descrição da alteração]\\
E3 & [Nome] & [Descrição da alteração]\\
\bottomrule
\end{tabular}
\end{table}

\subsection{Modelos pré-treinados e transferência de aprendizado}
% Instruções da especificação:
% Descreva as arquiteturas utilizadas no transfer learning, a origem dos pesos
% pré-treinados, a adaptação para as dez classes do CIFAR-10 e o pré-processamento.
% Descreva as estratégias de treinamento avaliadas, incluindo quais camadas foram
% congeladas e quais foram treinadas. Informe os hiperparâmetros e justifique as
% decisões tomadas sempre que possível.

[Descrição das arquiteturas pré-treinadas e das estratégias de treinamento.]

\subsection{Métricas e ambiente experimental}
% Instruções da especificação:
% Descreva as métricas de avaliação, o protocolo utilizado para os experimentos,
% o hardware e o software utilizados. A métrica principal deve ser accuracy;
% descreva também a geração e análise das matrizes de confusão. Informe, quando
% pertinente, parâmetros dos modelos, tempo de treinamento e tempo de inferência.

A acurácia é a métrica principal. As matrizes de confusão são calculadas para as
dez classes e permitem identificar erros recorrentes entre pares de classes.
Durante o treinamento, são registrados a perda e a acurácia de treinamento e
validação por época; essas séries são usadas para gerar as curvas apresentadas
na Seção de Resultados. Também são registrados o número total e o número de
parâmetros treináveis, o tempo total de treinamento e o tempo médio de
inferência por imagem. A medição de inferência inclui aquecimento do modelo e
sincronização quando executada em GPU.

As execuções registram as versões de Python e PyTorch, a versão de CUDA quando
disponível e o dispositivo utilizado. O hardware e as versões correspondentes
aos experimentos finais devem ser consolidados a partir dos arquivos de
resultado de cada execução: [informar os equipamentos e ambientes usados nos
experimentos finais].

% ============================================================
\section{Resultados e discussão}

\subsection{Resultados da CNN própria}
% Instruções da especificação:
% Apresente os resultados dos diferentes experimentos, incluindo, conforme pertinente,
% tabelas comparativas, curvas de treinamento e validação, matrizes de confusão,
% resultados por classe e análise de parâmetros/custo computacional.
% Não apresente apenas tabelas e gráficos: analise e discuta os resultados no texto.
% Procure responder:
% - Quais alterações na CNN própria produziram maior impacto?
% - Aumentar a capacidade da rede melhorou o desempenho?
% - As técnicas de regularização foram efetivas?
% - Como as alterações afetaram overfitting/underfitting e a generalização?
% Os resultados de teste não devem ser usados para escolher modelos ou hiperparâmetros.

[Apresente e discuta os resultados dos experimentos da CNN própria.]

% TABELA DE EXEMPLO: inclua apenas métricas e observações pertinentes aos
% experimentos realizados; os campos e as colunas podem ser alterados.
\begin{table}[H]
\centering
\caption{Exemplo de apresentação dos resultados de validação da CNN própria.}
\label{tab:resultados_cnn}
\small
\begin{tabular}{lrrrl}
\toprule
Modelo & Accuracy (\%) & Parâmetros & Tempo de treino & Observações\\
\midrule
Baseline & [ ] & [ ] & [ ] & [ ]\\
E1 & [ ] & [ ] & [ ] & [ ]\\
E2 & [ ] & [ ] & [ ] & [ ]\\
E3 & [ ] & [ ] & [ ] & [ ]\\
\bottomrule
\end{tabular}
\end{table}

% Exemplo de inserção de figura; substitua o caminho pelo arquivo real.
% \begin{figure}[H]
%     \centering
%     \includegraphics[width=0.85\linewidth]{figuras/curvas_cnn.pdf}
%     \caption{Curvas de treinamento e validação da CNN.}
%     \label{fig:curvas_cnn}
% \end{figure}

\subsection{Resultados de transferência de aprendizado}
% Instruções da especificação:
% Apresente e discuta os resultados das arquiteturas pré-treinadas e das estratégias
% avaliadas. Identifique as configurações selecionadas com base na validação.
% Discuta:
% - se os modelos pré-treinados apresentaram vantagem em relação à CNN própria;
% - se o congelamento ou descongelamento das camadas afetou o desempenho;
% - diferenças entre treinamento do zero e transfer learning, quando essa comparação
%   tiver sido realizada;
% - desempenho por classe, custo computacional e limitações observadas.

[Apresente e discuta os resultados dos modelos pré-treinados.]

% TABELA DE EXEMPLO: adapte as linhas às arquiteturas e estratégias efetivamente
% avaliadas; não é necessário manter todas as colunas sugeridas.
\begin{table}[H]
\centering
\caption{Exemplo de apresentação dos resultados de validação em transferência de aprendizado.}
\label{tab:transfer}
\small
\begin{tabular}{llrrr}
\toprule
Arquitetura & Estratégia & Accuracy val. (\%) & Parâmetros treináveis & Tempo\\
\midrule
Modelo 1 & [Estratégia 1] & [ ] & [ ] & [ ]\\
Modelo 1 & [Estratégia 2] & [ ] & [ ] & [ ]\\
Modelo 2 & [Estratégia 1] & [ ] & [ ] & [ ]\\
Modelo 2 & [Estratégia 2] & [ ] & [ ] & [ ]\\
\bottomrule
\end{tabular}
\end{table}

\subsection{Avaliação final no conjunto de teste}
% Instruções da especificação:
% Apresente a comparação global entre a melhor CNN própria e a melhor configuração
% de cada uma das duas arquiteturas pré-treinadas. Os resultados de teste devem
% corresponder às configurações selecionadas a partir da validação.
% Discuta:
% - quais classes foram mais difíceis de classificar e quais foram confundidas;
% - as diferenças entre as arquiteturas e estratégias;
% - a relação entre desempenho, número de parâmetros e custo computacional;
% - vantagens e limitações do transfer learning;
% - se há relação clara entre complexidade computacional e desempenho.
% Explique os resultados; não se limite a apresentar métricas e tabelas.

[Apresente a comparação final e discuta os resultados.]

% TABELA DE EXEMPLO: compare as configurações finais selecionadas com base na
% validação. Adapte as métricas e as linhas às abordagens realmente avaliadas.
\begin{table}[H]
\centering
\caption{Exemplo de comparação das configurações finais selecionadas no conjunto de teste.}
\label{tab:teste}
\small
\begin{tabular}{lrrr}
\toprule
Abordagem selecionada & Accuracy teste (\%) & Parâmetros & Tempo de inferência\\
\midrule
CNN própria & [ ] & [ ] & [ ]\\
Arquitetura 1 e estratégia & [ ] & [ ] & [ ]\\
Arquitetura 2 e estratégia & [ ] & [ ] & [ ]\\
\bottomrule
\end{tabular}
\end{table}

\subsection{Síntese e limitações}
% Instruções da especificação:
% Retome brevemente o que foi realizado no trabalho e apresente as principais
% conclusões obtidas a partir dos experimentos. Sintetize as diferenças observadas
% entre as arquiteturas e as conclusões sobre transfer learning. Descreva as
% dificuldades encontradas, limitações do trabalho e possíveis extensões futuras.

[Apresente a síntese dos achados e as limitações.]

% ============================================================
\section{Conclusão}
% Instruções da especificação:
% A conclusão deve conter:
% - uma retomada breve do que foi realizado no trabalho;
% - as principais conclusões obtidas a partir dos experimentos;
% - as principais diferenças observadas entre as arquiteturas;
% - conclusões sobre o uso de transfer learning;
% - dificuldades encontradas;
% - limitações do trabalho;
% - possíveis extensões ou experimentos futuros.
% Não introduza resultados novos nesta seção.

[Escreva a conclusão conforme as instruções acima.]

% ============================================================
\section*{Contribuição individual}
% Instrução da especificação:
% Liste as contribuições de cada integrante para o trabalho, explicitando
% o que cada um realizou.

\begin{itemize}[leftmargin=1.2cm,noitemsep,topsep=2pt]
    \item \textbf{[Nome do integrante 1]:} [Contribuições.]
    \item \textbf{[Nome do integrante 2]:} [Contribuições.]
    \item \textbf{[Nome do integrante 3]:} [Contribuições.]
    \item \textbf{[Nome do integrante 4, se aplicável]:} [Contribuições.]
\end{itemize}

% ============================================================
% REFERÊNCIAS
% Instruções:
% - Mantenha as referências em um arquivo externo chamado referencias.bib.
% - Cite as referências no texto com \cite{chave}.
% - Inclua no arquivo .bib apenas referências utilizadas no texto.
% - Compile com a sequência: pdfLaTeX -> BibTeX -> pdfLaTeX -> pdfLaTeX.
\bibliographystyle{plain}
\bibliography{referencias}

\end{document}
